\paragraph{\texorpdfstring{$2^L$}{2^L}-进全息像集的强分离几何与符号动力学}
沿用定理 \ref{thm:killo-godel-semidir-affine-2adic} 与推论 \ref{cor:killo-infinite-stream-2adic-holographic-point} 的记号。记
\[
A:=\mathrm{code}(\mathcal E)\subset\{0,1,\dots,M\},\qquad
q:=\card{A},\qquad
T:=T_2(E_{\min})\cong\ZZ_2^2,
\]
并令
\[
T_v:=\ZZ_2 v\subset T.
\]
由于 \(\ZZ_2\) 紧且 \(a\mapsto av\) 连续，\(T_v\) 为闭子群；又 \(T\) 无挠且 \(v\neq 0\)，故
\[
\iota_v:\ZZ_2\longrightarrow T_v,\qquad a\longmapsto av
\]
为拓扑群同构。下文以其传递的 \(2\)-进超度量
\[
d_v(av,bv):=|a-b|_2
\]
赋予 \(T_v\) 度量结构。通过编码将 \(\mathcal E^\NN\) 与 \(A^\NN\) 识别，并定义
\[
X(\mathbf a):=\sum_{t\ge 1}a_t B^{t-1}v,\qquad
\mathbf a=(a_t)_{t\ge 1}\in A^\NN,
\]
以及全息像集
\[
K_{\mathcal E}:=\{X(\mathbf a):\ \mathbf a\in A^\NN\}\subset T_v.
\]

\begin{theorem}[\texorpdfstring{$2^L$}{2^L}-进全息地址的精确柱分离]\label{thm:killo-2adic-holographic-exact-cylinder-separation}
进一步假设 \(v\) 为 \(T\) 中的原始向量，即
\[
v\notin BT.
\]
对任意 \(\mathbf a=(a_t)_{t\ge 1},\mathbf b=(b_t)_{t\ge 1}\in A^\NN\)，若其首次分歧位置为
\[
j:=\min\{t\ge 1:\ a_t\neq b_t\},
\]
则有
\[
X(\mathbf a)-X(\mathbf b)\in B^{j-1}T\setminus B^jT.
\]
等价地，对任意 \(n\ge 1\)，成立
\[
X(\mathbf a)\equiv X(\mathbf b)\pmod{B^nT}
\quad\Longleftrightarrow\quad
\mathrm{pref}_n(\mathbf a)=\mathrm{pref}_n(\mathbf b),
\]
其中 \(\mathrm{pref}_n\) 表示长度 \(n\) 前缀。
\leanverified{paper\_killo\_2adic\_holographic\_exact\_cylinder\_separation}
\end{theorem}
\begin{proof}
由 \(j\) 的定义，
\[
X(\mathbf a)-X(\mathbf b)
=
\sum_{t\ge j}(a_t-b_t)B^{t-1}v
=
B^{j-1}uv,
\]
其中
\[
u:=(a_j-b_j)+\sum_{t>j}(a_t-b_t)B^{t-j}\in \ZZ_2.
\]
由于 \(A\subset\{0,1,\dots,M\}\)、\(M<B\)，且 \(a_j\neq b_j\)，故
\[
u\equiv a_j-b_j\not\equiv 0\pmod B.
\]
因此 \(u\in\ZZ_2^\times\) 为单位元，从而 \(X(\mathbf a)-X(\mathbf b)\in B^{j-1}T\)。
若 \(X(\mathbf a)-X(\mathbf b)\in B^jT\)，则
\[
B^{j-1}uv\in B^jT=B^{j-1}(BT).
\]
由于 \(T\) 无挠，乘以 \(B^{j-1}\) 的映射在 \(T\) 上单射，故 \(uv\in BT\)。
再由 \(u\) 为单位元，乘以 \(u^{-1}\) 得 \(v\in BT\)，与假设矛盾。
故
\[
X(\mathbf a)-X(\mathbf b)\notin B^jT.
\]

若 \(\mathrm{pref}_n(\mathbf a)=\mathrm{pref}_n(\mathbf b)\)，则
\[
X(\mathbf a)-X(\mathbf b)=\sum_{t>n}(a_t-b_t)B^{t-1}v\in B^nT,
\]
故右端推出左端。
反之，若 \(X(\mathbf a)\equiv X(\mathbf b)\pmod{B^nT}\) 而前 \(n\) 位并不全同，则其首次分歧位置 \(j\le n\)，这与上面已证
\[
X(\mathbf a)-X(\mathbf b)\in B^{j-1}T\setminus B^jT
\]
矛盾。
故左端亦推出右端。
\end{proof}

\begin{theorem}[全息像集在 \texorpdfstring{$T_v$}{Tv} 上的强分离自相似结构与精确维数]\label{thm:killo-2adic-holographic-attractor-dimension}
\leanverified{paper\_killo\_2adic\_holographic\_attractor\_dimension}
对每个 \(a\in A\) 定义仿射压缩
\[
\Phi_a:T_v\to T_v,\qquad \Phi_a(x):=av+Bx.
\]
则成立：
\begin{enumerate}
  \item \(K_{\mathcal E}\) 是迭代函数系 \(\{\Phi_a\}_{a\in A}\) 的唯一非空紧吸引子，亦即
  \[
  K_{\mathcal E}=\bigcup_{a\in A}\Phi_a(K_{\mathcal E});
  \]
  \item 上述并为强分离并：若 \(a\neq a'\)，则
  \[
  \Phi_a(K_{\mathcal E})\cap \Phi_{a'}(K_{\mathcal E})=\varnothing;
  \]
  \item 以 \(d_v\) 计，
  \[
  \dim_H(K_{\mathcal E})=\dim_M(K_{\mathcal E})=\frac{\log q}{L\log 2};
  \]
  \item 若 \(\lambda_v\) 为 \(T_v\) 的 Haar 概率测度，则
  \[
  \lambda_v(K_{\mathcal E})=
  \begin{cases}
  0,& q<B,\\
  1,& q=B.
  \end{cases}
  \]
  并且在 \(q=B\) 时必有 \(A=\{0,1,\dots,B-1\}\)，从而 \(K_{\mathcal E}=T_v\)。
\end{enumerate}
\end{theorem}
\begin{proof}
对任意 \(\mathbf a=(a_t)_{t\ge 1}\in A^\NN\)，若记 \(\sigma(\mathbf a)=(a_{t+1})_{t\ge 1}\)，则
\[
X(\mathbf a)=a_1v+B\,X(\sigma\mathbf a).
\]
因此
\[
K_{\mathcal E}=\bigcup_{a\in A}(av+BK_{\mathcal E})=\bigcup_{a\in A}\Phi_a(K_{\mathcal E}).
\]
在 \(\iota_v\) 坐标下，
\[
\iota_v^{-1}\circ \Phi_a\circ \iota_v(y)=a+By,
\]
故每个 \(\Phi_a\) 都是相似比 \(2^{-L}\) 的严格压缩；由 Hutchinson 定理得 \(K_{\mathcal E}\) 为唯一非空紧吸引子。

再证强分离。若 \(x\in \Phi_a(K_{\mathcal E})\cap \Phi_{a'}(K_{\mathcal E})\)，则
\[
x-av\in BT_v,\qquad x-a'v\in BT_v,
\]
从而
\[
(a-a')v\in BT_v.
\]
经 \(\iota_v^{-1}\) 拉回得 \(a-a'\in B\ZZ_2\cap \ZZ=B\ZZ\)。由于 \(a,a'\in\{0,1,\dots,B-1\}\)，故 \(|a-a'|<B\)，只能有 \(a-a'=0\)，即 \(a=a'\)。

对长度 \(n\) 的字 \(w=(a_1,\dots,a_n)\in A^n\)，记
\[
K_w:=\Phi_{a_1}\circ\cdots\circ\Phi_{a_n}(K_{\mathcal E}).
\]
若再记
\[
s_w:=\sum_{j=1}^{n}a_jB^{j-1},
\]
则
\[
K_w\subset s_wv+B^nT_v.
\]
若 \(w\neq w'\)，则由前述逐位消去论证可知 \(s_w-s_{w'}\notin B^n\ZZ\)，从而相应的 \(B^nT_v\)-陪集不同。
每个 \(K_w\) 都是 \(K_{\mathcal E}\) 的相似拷贝，且
\[
\operatorname{diam}_{d_v}(K_w)=2^{-Ln}\operatorname{diam}_{d_v}(K_{\mathcal E}).
\]
又由强分离，\(\{K_w\}_{w\in A^n}\) 两两不交，个数恰为 \(q^n\)。因此经典 Moran 公式适用，临界指数 \(s\) 满足
\[
q\cdot 2^{-Ls}=1,
\]
故
\[
s=\frac{\log q}{L\log 2}.
\]
于是 Hausdorff 维数与 Minkowski 维数同为该值。

最后证明 Haar 测度结论。每个 \(K_w\) 都包含于某个 \(B^nT_v\)-陪集，而不同 \(w\in A^n\) 对应不同陪集；\(T_v/B^nT_v\) 恰有 \(B^n\) 个陪集，每个陪集的 Haar 测度均为 \(B^{-n}\)。故对每个 \(n\),
\[
\lambda_v(K_{\mathcal E})\le q^n B^{-n}.
\]
若 \(q<B\)，令 \(n\to\infty\) 即得 \(\lambda_v(K_{\mathcal E})=0\)。

若 \(q=B\)，则 \(A\subseteq\{0,1,\dots,B-1\}\) 且基数为 \(B\)，故必有 \(A=\{0,1,\dots,B-1\}\)。此时经 \(\iota_v^{-1}\) 识别，
\[
\iota_v^{-1}(K_{\mathcal E})
=
\left\{\sum_{t\ge 1}a_tB^{t-1}:\ a_t\in\{0,1,\dots,B-1\}\right\}
=\ZZ_2,
\]
即 \(K_{\mathcal E}=T_v\)，从而 \(\lambda_v(K_{\mathcal E})=1\)。证毕。
\end{proof}

\begin{corollary}[原始地址线在环境 \texorpdfstring{$T$}{T} 中的精确维数与零测性]\label{cor:killo-2adic-holographic-ambient-dimension-zero}
在定理 \ref{thm:killo-2adic-holographic-exact-cylinder-separation} 的原始向量假设下，\(q=\card{A}=\card{\mathcal E}\)，并且把 \(K_{\mathcal E}\) 视为环境空间 \(T\cong \ZZ_2^2\) 的子集时，仍有
\[
\dim_H(K_{\mathcal E})=\dim_M(K_{\mathcal E})=\frac{\log q}{L\log 2}=\frac{\log \card{\mathcal E}}{\log B}.
\]
此外，由 \(A\subset\{0,1,\dots,M\}\) 与 \(B>M\) 可知 \(q\le B\)，故
\[
\dim_H(K_{\mathcal E})\le 1<2=\dim_H(T).
\]
因此 \(K_{\mathcal E}\) 在 \(T\) 的 Haar 测度下必为零测集。
\end{corollary}
\begin{proof}
由编码单射性，\(\card{A}=\card{\mathcal E}\)。
将 \(v\) 唯一写成
\[
v=2^r u,
\]
其中 \(r\ge 0\)，且 \(u\in T\setminus 2T\) 为 \(\ZZ_2\)-模意义下的原始向量。
于是可取 \(w\in T\) 使
\[
T=\ZZ_2u\oplus \ZZ_2w.
\]
在该基下定义环境超度量
\[
d_T(x_1u+y_1w,\ x_2u+y_2w):=\max\{|x_1-x_2|_2,\ |y_1-y_2|_2\}.
\]
则 \(d_T\) 与 \(T\) 的标准 \(2\)-进拓扑相容。
而对任意 \(a,b\in \ZZ_2\)，有
\[
d_T(av,bv)=d_T(2^rau,2^rbu)=2^{-r}|a-b|_2.
\]
因此 \(d_T|_{T_v}\) 与 \(d_v\) 只差一个常数因子 \(2^{-r}\)，从而二者双 Lipschitz 等价。
因此 \(K_{\mathcal E}\subset T_v\subset T\) 的 Hausdorff 维数与 Minkowski 维数，在 \(T_v\) 内与在 \(T\) 内计算时一致。
由定理 \ref{thm:killo-2adic-holographic-attractor-dimension}，得到
\[
\dim_H(K_{\mathcal E})=\dim_M(K_{\mathcal E})=\frac{\log q}{L\log 2}.
\]

再由 \(A\subset\{0,1,\dots,M\}\) 与 \(B>M\)，得
\[
q=\card{A}\le M+1\le B.
\]
故
\[
\frac{\log q}{\log B}\le 1.
\]
另一方面，\(T\cong \ZZ_2^2\) 在上述乘积超度量下的 Hausdorff 维数为 \(2\)，于是
\[
\dim_H(K_{\mathcal E})\le 1<2=\dim_H(T).
\]
由标准维数判别，\(K_{\mathcal E}\) 在 \(T\) 的 Haar 测度下为零测集。证毕。
\end{proof}
\leanverified{paper\_killo\_2adic\_holographic\_ambient\_dimension\_zero}

\begin{proposition}[前缀分辨率与 \texorpdfstring{$B^nT_v$}{B\^{}nTv}-陪集的严格分类]\label{prop:killo-2adic-holographic-prefix-classification}
\leanverified{paper\_killo\_2adic\_holographic\_prefix\_classification}
对每个 \(n\ge 0\)，记
\[
\rho_n:T_v\to T_v/B^nT_v
\]
为自然投影。对每个 \(n\ge 1\)，定义前缀映射
\[
\mathrm{pref}_n:A^\NN\to A^n,\qquad
\mathrm{pref}_n(\mathbf a)=(a_1,\dots,a_n),
\]
则对任意 \(\mathbf a,\mathbf b\in A^\NN\)，成立等价关系
\[
\rho_n(X(\mathbf a))=\rho_n(X(\mathbf b))
\quad\Longleftrightarrow\quad
\mathrm{pref}_n(\mathbf a)=\mathrm{pref}_n(\mathbf b).
\]
因此 \(\rho_n(K_{\mathcal E})\) 恰有 \(q^n\) 个点。若 \(q\ge 2\)，则对每个 \(n\ge 1\)，商映射 \(\rho_{n-1}|_{K_{\mathcal E}}\) 不能区分全部长度 \(n\) 的前缀。
\end{proposition}
\begin{proof}
若 \(\mathrm{pref}_n(\mathbf a)=\mathrm{pref}_n(\mathbf b)\)，则
\[
X(\mathbf a)-X(\mathbf b)=\sum_{t>n}(a_t-b_t)B^{t-1}v\in B^nT_v,
\]
故 \(\rho_n(X(\mathbf a))=\rho_n(X(\mathbf b))\)。

反之，若 \(\rho_n(X(\mathbf a))=\rho_n(X(\mathbf b))\)，则
\[
X(\mathbf a)-X(\mathbf b)\in B^nT_v.
\]
经 \(\iota_v^{-1}\) 拉回得
\[
\sum_{t=1}^{n}(a_t-b_t)B^{t-1}\in B^n\ZZ_2.
\]
模 \(B\) 约化可知 \(a_1-b_1\equiv 0\pmod B\)。由于 \(a_1,b_1\in\{0,1,\dots,B-1\}\)，只能有 \(a_1=b_1\)。减去首项并除以 \(B\)，得到
\[
\sum_{t=2}^{n}(a_t-b_t)B^{t-2}\in B^{n-1}\ZZ_2.
\]
重复此论证即得 \(a_t=b_t\)（\(1\le t\le n\)），故 \(\mathrm{pref}_n(\mathbf a)=\mathrm{pref}_n(\mathbf b)\)。

于是长度 \(n\) 的前缀与 \(K_{\mathcal E}\) 在 \(T_v/B^nT_v\) 中的有效像一一对应，故 \(\card{\rho_n(K_{\mathcal E})}=q^n\)。

若 \(q\ge 2\)，取 \(a\neq b\in A\) 及任意固定尾部 \(\mathbf c\in A^\NN\)。则
\[
\mathbf a:=(c_1,\dots,c_{n-1},a,c_1,c_2,\dots),\qquad
\mathbf b:=(c_1,\dots,c_{n-1},b,c_1,c_2,\dots)
\]
具有相同的长度 \(n-1\) 前缀，从而
\[
X(\mathbf a)-X(\mathbf b)\in B^{n-1}T_v,
\]
即 \(\rho_{n-1}(X(\mathbf a))=\rho_{n-1}(X(\mathbf b))\)，但两者长度 \(n\) 前缀不同。故 \(\rho_{n-1}|_{K_{\mathcal E}}\) 不能区分全部长度 \(n\) 的前缀。证毕。
\end{proof}

\begin{theorem}[全息地址像的 cylinder 熵维数与分形维数完全一致]\label{thm:killo-2adic-holographic-cylinder-entropy-dimension}
沿用定理 \ref{thm:killo-2adic-holographic-attractor-dimension} 与命题 \ref{prop:killo-2adic-holographic-prefix-classification} 的记号。则极限
\[
\lim_{n\to\infty}\frac{\log \card{\rho_n(K_{\mathcal E})}}{nL\log 2}
\]
存在，并满足
\[
\dim_H(K_{\mathcal E})
=
\dim_M(K_{\mathcal E})
=
\lim_{n\to\infty}\frac{\log \card{\rho_n(K_{\mathcal E})}}{nL\log 2}
=
\frac{\log q}{L\log 2}.
\]
因此，Hausdorff 维数、Minkowski 维数与由 cylinder 计数定义的熵维数在该 \(2^L\)-进地址像上完全重合；特别地，\(K_{\mathcal E}\) 在每一层都由 \(q\) 个相似比 \(2^{-L}\) 的两两不交分支严格均匀生成。
\leanverified{paper\_killo\_2adic\_holographic\_cylinder\_entropy\_dimension}
\end{theorem}
\begin{proof}
定理 \ref{thm:killo-2adic-holographic-attractor-dimension} 已证明
\[
\dim_H(K_{\mathcal E})=\dim_M(K_{\mathcal E})=\frac{\log q}{L\log 2}.
\]
另一方面，命题 \ref{prop:killo-2adic-holographic-prefix-classification} 给出对每个 \(n\ge 0\),
\[
\card{\rho_n(K_{\mathcal E})}=q^n.
\]
故
\[
\frac{\log \card{\rho_n(K_{\mathcal E})}}{nL\log 2}
=
\frac{\log q^n}{nL\log 2}
=
\frac{\log q}{L\log 2}
\qquad(n\ge 1),
\]
从而所示极限存在并等于同一个值。

最后，定理 \ref{thm:killo-2adic-holographic-attractor-dimension} 已给出
\[
K_{\mathcal E}=\bigsqcup_{a\in A}\Phi_a(K_{\mathcal E}),
\qquad
\Phi_a(x)=av+Bx,
\]
且每个 \(\Phi_a\) 都是相似比 \(2^{-L}\) 的严格压缩。由此得到最后的均匀分支表述。证毕。
\end{proof}

\begin{theorem}[全息地址像的 cylinder \texorpdfstring{$\zeta$}{zeta} 函数与格型复维数]\label{thm:killo-2adic-holographic-cylinder-zeta-complex-dimensions}
沿用定理 \ref{thm:killo-2adic-holographic-cylinder-entropy-dimension} 的记号，并定义 cylinder \(\zeta\) 函数
\[
\zeta_{K_{\mathcal E}}^{\mathrm{cyl}}(s)
:=
\sum_{n\ge 0}\card{\rho_n(K_{\mathcal E})}\,2^{-Lns}
=
\sum_{n\ge 0}q^n\,2^{-Lns}.
\]
则在半平面
\[
\Re s>\frac{\log q}{L\log 2}
\]
上有闭式
\[
\zeta_{K_{\mathcal E}}^{\mathrm{cyl}}(s)=\frac{1}{1-q\,2^{-Ls}}.
\]
其收敛临界线恰为
\[
\Re s=\frac{\log q}{L\log 2},
\]
并且全部极点构成一条竖直格：
\[
\mathcal D_{K_{\mathcal E}}
=
\left\{
\frac{\log q}{L\log 2}
\;+\;
\frac{2\pi i k}{L\log 2}:\ k\in\ZZ
\right\}.
\]
\end{theorem}
\begin{proof}
由命题 \ref{prop:killo-2adic-holographic-prefix-classification}，
\[
\card{\rho_n(K_{\mathcal E})}=q^n
\qquad(n\ge 0).
\]
代入定义即可得到
\[
\zeta_{K_{\mathcal E}}^{\mathrm{cyl}}(s)
=
\sum_{n\ge 0}\bigl(q\,2^{-Ls}\bigr)^n.
\]
故当且仅当
\[
\bigl|q\,2^{-Ls}\bigr|=q\,2^{-L\Re s}<1
\]
时级数绝对收敛，也即
\[
\Re s>\frac{\log q}{L\log 2}.
\]
在此区域内，几何级数求和给出
\[
\zeta_{K_{\mathcal E}}^{\mathrm{cyl}}(s)=\frac{1}{1-q\,2^{-Ls}}.
\]

极点由方程
\[
q\,2^{-Ls}=1
\]
确定。写成指数形式即
\[
-Ls\log 2=-\log q+2\pi i k,
\qquad k\in\ZZ,
\]
从而
\[
s=\frac{\log q}{L\log 2}-\frac{2\pi i k}{L\log 2}.
\]
当 \(k\) 遍历 \(\ZZ\) 时，这一集合等同于
\[
\left\{
\frac{\log q}{L\log 2}
\;+\;
\frac{2\pi i k}{L\log 2}:\ k\in\ZZ
\right\}.
\]
故全部复维数恰构成所示竖直格。证毕。
\end{proof}

\begin{theorem}[Bernoulli 全息推前测度的精确局部维数]\label{thm:killo-2adic-holographic-bernoulli-exact-dimension}
取 \(A\) 上任意严格正的概率向量
\[
p=(p_a)_{a\in A},\qquad p_a>0,\qquad \sum_{a\in A}p_a=1,
\]
并令 \(\mathbb P_p:=p^\NN\) 为 \(A^\NN\) 上的 Bernoulli 乘积测度。定义全息推前测度
\[
\mu_p:=X_*\mathbb P_p
\]
于 \(K_{\mathcal E}\) 上。则 \(\mu_p\) 为 exact-dimensional 测度，并且对 \(\mu_p\)-几乎处处的 \(x\in K_{\mathcal E}\)，局部维数
\[
\dim_{\mathrm{loc}}(\mu_p,x)
:=
\lim_{r\downarrow 0}
\frac{\log \mu_p(B_{d_v}(x,r))}{\log r}
\]
存在且满足
\[
\dim_{\mathrm{loc}}(\mu_p,x)=\frac{H(p)}{L\log 2},
\qquad
H(p):=-\sum_{a\in A}p_a\log p_a.
\]
特别地，
\[
\dim_H(\mu_p)=\frac{H(p)}{L\log 2}.
\]
当 \(p\) 为 \(A\) 上的均匀分布时，
\[
\dim_H(\mu_p)=\frac{\log q}{L\log 2}=\dim_H(K_{\mathcal E}),
\]
故均匀 Bernoulli 推前为 \(K_{\mathcal E}\) 的满维测度。
\end{theorem}
\begin{proof}
令 \(x=X(\mathbf a)\)，其中 \(\mathbf a=(a_t)_{t\ge 1}\in A^\NN\)。由命题 \ref{prop:killo-2adic-holographic-prefix-classification}，半径 \(2^{-Ln}\) 的 \(d_v\)-球在 \(K_{\mathcal E}\) 上恰对应长度 \(n\) 前缀类，故
\[
\mu_p\!\bigl(B_{d_v}(x,2^{-Ln})\bigr)
=
\mathbb P_p\!\bigl(\mathrm{pref}_n^{-1}(a_1,\dots,a_n)\bigr)
=
\prod_{t=1}^{n}p_{a_t}.
\]
于是
\[
\frac{\log \mu_p(B_{d_v}(x,2^{-Ln}))}{\log 2^{-Ln}}
=
\frac{\sum_{t=1}^{n}\log p_{a_t}}{-Ln\log 2}.
\]
由大数定律，\(\mathbb P_p\)-几乎处处有
\[
\frac1n\sum_{t=1}^{n}\log p_{a_t}\longrightarrow \sum_{a\in A}p_a\log p_a=-H(p).
\]
因此对 \(\mu_p\)-几乎处处的 \(x\in K_{\mathcal E}\),
\[
\dim_{\mathrm{loc}}(\mu_p,x)=\frac{H(p)}{L\log 2}.
\]
这表明 \(\mu_p\) 为 exact-dimensional，且其 Hausdorff 维数等于该常数。均匀分布时 \(H(p)=\log q\)，遂得最后结论。证毕。
\end{proof}

\begin{theorem}[全息移位动力学的熵、周期点与 Artin--Mazur \texorpdfstring{$\zeta$}{zeta} 函数]\label{thm:killo-2adic-holographic-shift-dynamics}
设 \(\sigma:A^\NN\to A^\NN\) 为左移
\[
\sigma(a_1,a_2,a_3,\dots)=(a_2,a_3,\dots).
\]
定义
\[
T_{\mathcal E}:K_{\mathcal E}\to K_{\mathcal E},\qquad
T_{\mathcal E}(X(\mathbf a)):=X(\sigma\mathbf a).
\]
则：
\begin{enumerate}
  \item \(T_{\mathcal E}\) 为良定义连续映射，且 \(X:A^\NN\to K_{\mathcal E}\) 为把 \(\sigma\) 共轭到 \(T_{\mathcal E}\) 的同胚；
  \item 对每个 \(a\in A\)，若记 \(K_a:=\Phi_a(K_{\mathcal E})\)，则在分支 \(K_a\) 上有显式表达
  \[
  T_{\mathcal E}(x)=B^{-1}(x-av),\qquad x\in K_a;
  \]
  \item 拓扑熵满足
  \[
  h_{\mathrm{top}}(T_{\mathcal E})=\log q;
  \]
  \item 对每个 \(n\ge 1\)，周期点计数为
  \[
  \card{\mathrm{Fix}(T_{\mathcal E}^n)}=q^n;
  \]
  \item Artin--Mazur \(\zeta\) 函数闭式为
  \[
  \zeta_{T_{\mathcal E}}(z)
  :=
  \exp\!\left(\sum_{n\ge 1}\frac{\card{\mathrm{Fix}(T_{\mathcal E}^n)}}{n}z^n\right)
  =
  \frac{1}{1-qz}.
  \]
\end{enumerate}
\end{theorem}

\leanverified{paper\_killo\_2adic\_holographic\_shift\_dynamics}

\begin{proof}
推论 \ref{cor:killo-infinite-stream-2adic-holographic-point} 已给出 \(X\) 的单射性；而 \(A^\NN\) 紧、\(K_{\mathcal E}\subset T_v\) Hausdorff，故连续双射
\[
X:A^\NN\to K_{\mathcal E}
\]
为同胚。于是 \(T_{\mathcal E}:=X\circ \sigma\circ X^{-1}\) 良定义且连续，并自动满足
\[
T_{\mathcal E}\circ X=X\circ \sigma.
\]

另一方面，若 \(x\in K_a\)，则由定理 \ref{thm:killo-2adic-holographic-attractor-dimension} 的强分离分解可唯一写成
\[
x=av+By,\qquad y\in K_{\mathcal E}.
\]
故
\[
T_{\mathcal E}(x)=y=B^{-1}(x-av).
\]
这给出分支上的显式仿射表达。

由于 \((K_{\mathcal E},T_{\mathcal E})\) 与 \(q\) 符号的一边满移位 \((A^\NN,\sigma)\) 拓扑共轭，标准符号动力学给出
\[
h_{\mathrm{top}}(T_{\mathcal E})=h_{\mathrm{top}}(\sigma)=\log q.
\]
同理，\(\sigma^n\) 的不动点恰为长度 \(n\) 周期序列，其数目为 \(q^n\)，故
\[
\card{\mathrm{Fix}(T_{\mathcal E}^n)}=q^n.
\]
代入 Artin--Mazur \(\zeta\) 函数定义，得到
\[
\zeta_{T_{\mathcal E}}(z)
=
\exp\!\left(\sum_{n\ge 1}\frac{q^n}{n}z^n\right)
=
\exp\!\bigl(-\log(1-qz)\bigr)
=
\frac{1}{1-qz}.
\]
证毕。
\end{proof}

\begin{remark}[地址化链条的几何--动力学闭包]\label{rem:killo-2adic-holographic-closure}
定理 \ref{thm:killo-godel-semidir-affine-2adic} 与推论 \ref{cor:killo-infinite-stream-2adic-holographic-point} 中的 \(2^L\)-进仿射表示，在 \(T_v\) 上闭合为
\[
\text{仿射注入}
\Longrightarrow
\text{强分离自相似像集}
\Longrightarrow
\text{精确维数与局部维数}
\Longrightarrow
\text{前缀分辨率最优性}
\Longrightarrow
\text{移位共轭与 }\zeta\text{ 闭式}.
\]
这一闭包与正文中 Stokes--边界--能量--相位接口互补：前者给出地址空间的超度量几何与符号动力学公式，后者给出边界可审计性与误差控制的解析读出。
\end{remark}
